% GENERATED FILE. Edit canonical lessons, then run npm run generate:books.
% canonical-source-hash: fnv1a64:a7db5f9874d97b7f
% canonical-lessons: RU-C156-bogatyy, RU-C156-bednyy, RU-C156-chistyy, RU-C156-gryaznyy, RU-C156-pustoy

\chapter{Rich, Poor, Clean, Dirty, Empty}
\label{ch:ru-rich-poor-clean-dirty-empty}
\hlchaptermodality{156}

\begin{chapteropening}
\textbf{By the end of this chapter:} \emph{I can say rich, poor, clean, dirty and empty.}

Complete the last lesson of chapter 156: I can say rich, poor, clean, dirty and empty.
\end{chapteropening}

\section[\texorpdfstring{\ru{богатый} (bogátyy)}{богатый (bogátyy)}]{\ru{богатый} (bogátyy) — rich}

\label{lesson:RU-C156-bogatyy}



\begin{quote}
Before the new one: say the Russian for lazy, then the Russian for polite.
\end{quote}

\subsection*{You'll want to know: \ru{богатый}}
\textbf{\ru{богатый}} (\emph{bogátyy}) — \textquotedblleft{}rich\textquotedblright{}.

\begin{cousinweb}
\textbf{\ru{богатый}} comes from Proto-Slavic \emph{*bogŭ}: Vasmer gives both readings, \textquotedblleft{}god\textquotedblright{}, so \textquotedblleft{}favoured by the gods\textquotedblright{} like Latin \emph{fortūnātus}, and \textquotedblleft{}a share, wealth\textquotedblright{}, as in \emph{\ru{убогий}}, \textquotedblleft{}poor\textquotedblright{}, \textquotedblleft{}without a share\textquotedblright{}.
\end{cousinweb}

\begin{grammarlens}[title={What you've built}]
One more word for this chapter.
\end{grammarlens}

\subsection*{Your turn}
\begin{itemize}
\raggedright
  \item \emph{Say it:} \emph{bogátyy}
  \item \emph{Say it:} \emph{bogátyy}, once more
  \item \emph{Say it:} say \emph{bogátyy}
  \item \emph{Recall:} say the Russian for lazy, then the Russian for polite, then say \emph{bogátyy} again
\end{itemize}

\begin{tcolorbox}[breakable,colback=teal!4,colframe=teal!35!black,title={Before you move on}]
What does \ru{богатый} mean? (\textquotedblleft{}Rich\textquotedblright{}.) Say it once more.
\end{tcolorbox}
\section[\texorpdfstring{\ru{бедный} (bédnyy)}{бедный (bédnyy)}]{\ru{бедный} (bédnyy) — poor}

\label{lesson:RU-C156-bednyy}



\begin{quote}
Before the new one: say the Russian for polite, then the Russian for rich.
\end{quote}

\subsection*{You'll want to know: \ru{бедный}}
\textbf{\ru{бедный}} (\emph{bédnyy}) — \textquotedblleft{}poor\textquotedblright{}.

\begin{cousinweb}
\textbf{\ru{бедный}} is \emph{\ru{беда}}, \textquotedblleft{}misfortune\textquotedblright{}, with \emph{-\ru{ный}}: \textquotedblleft{}unfortunate\textquotedblright{}, then \textquotedblleft{}poor\textquotedblright{}. Vasmer compares Lithuanian \emph{bėdà}, \textquotedblleft{}trouble\textquotedblright{}, and Latvian \emph{będa}, \textquotedblleft{}sorrow\textquotedblright{}.
\end{cousinweb}

\begin{grammarlens}[title={What you've built}]
One more word for this chapter.
\end{grammarlens}

\subsection*{Your turn}
\begin{itemize}
\raggedright
  \item \emph{Say it:} \emph{bédnyy}
  \item \emph{Say it:} \emph{bédnyy}, once more
  \item \emph{Say it:} say \emph{bédnyy}
  \item \emph{Recall:} say the Russian for polite, then the Russian for rich, then say \emph{bédnyy} again
\end{itemize}

\begin{tcolorbox}[breakable,colback=teal!4,colframe=teal!35!black,title={Before you move on}]
What does \ru{бедный} mean? (\textquotedblleft{}Poor\textquotedblright{}.) Say it once more.
\end{tcolorbox}
\section[\texorpdfstring{\ru{чистый} (chístyy)}{чистый (chístyy)}]{\ru{чистый} (chístyy) — clean}

\label{lesson:RU-C156-chistyy}



\begin{quote}
Before the new one: say the Russian for rich, then the Russian for poor.
\end{quote}

\subsection*{You'll want to know: \ru{чистый}}
\textbf{\ru{чистый}} (\emph{chístyy}) — \textquotedblleft{}clean\textquotedblright{}.

\begin{cousinweb}
\textbf{\ru{чистый}} $\leftarrow$ Proto-Slavic \emph{*čistŭ} — inherited. Vasmer relates it to Old Prussian \emph{skijstan}, \textquotedblleft{}clean\textquotedblright{}, and Lithuanian \emph{skýstas}, \textquotedblleft{}liquid\textquotedblright{}, and perhaps to \emph{\ru{цедить}}, \textquotedblleft{}to strain\textquotedblright{}.
\end{cousinweb}

\begin{grammarlens}[title={What you've built}]
One more word for this chapter.
\end{grammarlens}

\subsection*{Your turn}
\begin{itemize}
\raggedright
  \item \emph{Say it:} \emph{chístyy}
  \item \emph{Say it:} \emph{chístyy}, once more
  \item \emph{Say it:} say \emph{chístyy}
  \item \emph{Recall:} say the Russian for rich, then the Russian for poor, then say \emph{chístyy} again
\end{itemize}

\begin{tcolorbox}[breakable,colback=teal!4,colframe=teal!35!black,title={Before you move on}]
What does \ru{чистый} mean? (\textquotedblleft{}Clean\textquotedblright{}.) Say it once more.
\end{tcolorbox}
\section[\texorpdfstring{\ru{грязный} (gryáznyy)}{грязный (gryáznyy)}]{\ru{грязный} (gryáznyy) — dirty}

\label{lesson:RU-C156-gryaznyy}



\begin{quote}
Before the new one: say the Russian for poor, then the Russian for clean.
\end{quote}

\subsection*{You'll want to know: \ru{грязный}}
\textbf{\ru{грязный}} (\emph{gryáznyy}) — \textquotedblleft{}dirty\textquotedblright{}.

\begin{cousinweb}
\textbf{\ru{грязный}} is \emph{\ru{грязь}}, \textquotedblleft{}mud\textquotedblright{}, with \emph{-\ru{ный}}. Vasmer puts \emph{\ru{грязь}} with \emph{\ru{грязнуть}}, \textquotedblleft{}to get stuck, to sink in\textquotedblright{}: mud is what you sink into.
\end{cousinweb}

\begin{grammarlens}[title={What you've built}]
One more word for this chapter.
\end{grammarlens}

\subsection*{Your turn}
\begin{itemize}
\raggedright
  \item \emph{Say it:} \emph{gryáznyy}
  \item \emph{Say it:} \emph{gryáznyy}, once more
  \item \emph{Say it:} say \emph{gryáznyy}
  \item \emph{Recall:} say the Russian for poor, then the Russian for clean, then say \emph{gryáznyy} again
\end{itemize}

\begin{tcolorbox}[breakable,colback=teal!4,colframe=teal!35!black,title={Before you move on}]
What does \ru{грязный} mean? (\textquotedblleft{}Dirty\textquotedblright{}.) Say it once more.
\end{tcolorbox}
\section[\texorpdfstring{\ru{пустой} (pustóy)}{пустой (pustóy)}]{\ru{пустой} (pustóy) — empty}

\label{lesson:RU-C156-pustoy}



\begin{quote}
Before the new one: say the Russian for clean, then the Russian for dirty.
\end{quote}

\subsection*{You'll want to know: \ru{пустой}}
\textbf{\ru{пустой}} (\emph{pustóy}) — \textquotedblleft{}empty\textquotedblright{}.

\begin{cousinweb}
\textbf{\ru{пустой}} $\leftarrow$ Proto-Slavic \emph{*pustŭ} — inherited. Vasmer relates it to Old Prussian \emph{pausto}, \textquotedblleft{}wild\textquotedblright{}, and \emph{paustre}, \textquotedblleft{}a wilderness\textquotedblright{}.
\end{cousinweb}

\begin{grammarlens}[title={What you've built}]
One more word for this chapter.
\end{grammarlens}

\subsection*{Your turn}
\begin{itemize}
\raggedright
  \item \emph{Say it:} \emph{pustóy}
  \item \emph{Say it:} \emph{pustóy}, once more
  \item \emph{Say it:} say \emph{pustóy}
  \item \emph{Recall:} say the Russian for clean, then the Russian for dirty, then say \emph{pustóy} again
\end{itemize}

\begin{tcolorbox}[breakable,colback=teal!4,colframe=teal!35!black,title={Before you move on}]
What does \ru{пустой} mean? (\textquotedblleft{}Empty\textquotedblright{}.) Say it once more.
\end{tcolorbox}
